\documentclass[preview]{standalone}

\usepackage{amsmath}
\usepackage{amssymb}
\usepackage{stellar}
\usepackage{definitions}
\usepackage{bettelini}

\begin{document}

\id{sequences-functions-boundedness}
\genpage

\section{Boundedness}

\begin{snippetdefinition}{function-sequence-boundedness-definition}{Bounded sequence of functions}
    Let \(S\subseteq\realnumbers\) and let
    \(\{f_n\}_{n\in\naturalnumbers}\) be a \sequence of \function[functions]
    \(f_n\colon S\fromto\realnumbers\). The sequence is \emph{bounded} on \(S\) if
    \[
        \forall n\in\naturalnumbers,\ \exists M_n>0\ \suchthat
        \forall x\in S,\quad |f_n(x)|\leq M_n.
    \]
\end{snippetdefinition}

\begin{snippetdefinition}{function-sequence-uniform-boundedness-definition}{Uniformly bounded sequence of functions}
    Let \(S\subseteq\realnumbers\) and let
    \(\{f_n\}_{n\in\naturalnumbers}\) be a \sequence of \function[functions]
    \(f_n\colon S\fromto\realnumbers\). The sequence is \emph{uniformly bounded}
    on \(S\) if
    \[
        \exists M>0\ \suchthat\ \forall n\in\naturalnumbers,\ \forall x\in S,
        \quad |f_n(x)|\leq M.
    \]
\end{snippetdefinition}

\begin{plainhtml}
    Uniform boundedness implies that every function in the sequence is bounded; the converse need not hold.
\end{plainhtml}

\begin{snippetproposition}{uniform-limit-boundedness-proposition}{Boundedness of a uniform limit}
    Let \(S\subseteq\realnumbers\) and let
    \(\{f_n\}_{n\in\naturalnumbers}\) be a \sequence of bounded
    \function[functions] \(f_n\colon S\fromto\realnumbers\).
    If \(f_n\) converges uniformly on \(S\) to
    \(f\colon S\fromto\realnumbers\), then \(f\) is bounded on \(S\) and
    \(\{f_n\}_{n\in\naturalnumbers}\) is uniformly bounded on \(S\).
\end{snippetproposition}

\begin{snippetproof}{uniform-limit-boundedness-proposition-proof}{uniform-limit-boundedness-proposition}{Boundedness of a uniform limit}
    By uniform convergence with \(\varepsilon=1\), there exists
    \(N\in\naturalnumbers\) such that
    \[
        |f_n(x)-f(x)|<1,
        \qquad \forall n>N,\ \forall x\in S.
    \]
    Fix \(n_0>N\). Since \(f_{n_0}\) is bounded, there exists
    \(M_{n_0}>0\) such that \(|f_{n_0}(x)|\leq M_{n_0}\) for every
    \(x\in S\). Hence
    \[
        |f(x)|\leq |f(x)-f_{n_0}(x)|+|f_{n_0}(x)|<1+M_{n_0},
    \]
    so \(f\) is bounded. Choose \(M>0\) such that \(|f(x)|\leq M\) on
    \(S\). For every \(n>N\),
    \[
        |f_n(x)|\leq |f_n(x)-f(x)|+|f(x)|<1+M.
    \]
    For the finitely many indices \(1,\ldots,N\), choose bounds
    \(M_1,\ldots,M_N\). Then
    \[
        M'=\max\{M_1,\ldots,M_N,1+M\}
    \]
    bounds \(|f_n(x)|\) for every \(n\in\naturalnumbers\) and \(x\in S\).
\end{snippetproof}

\section{Counterexamples}

\begin{snippetexample}{uniform-boundedness-without-pointwise-convergence-example}{Uniform boundedness does not imply pointwise convergence}
    On \(S=[0,2\pi]\), let \(f_n(x)=\sin(nx)\). The sequence is uniformly
    bounded by \(1\), but it does not converge pointwise. Indeed, at
    \(x=\frac{\pi}{2}\),
    \[
        \sin\left(n\frac{\pi}{2}\right)=
        \begin{cases}
            1 & n=1+4k,\\
            0 & n=2k,\\
            -1 & n=3+4k.
        \end{cases}
    \]
\end{snippetexample}

\begin{snippetexample}{uniform-boundedness-without-pointwise-convergent-subsequence-example}{Uniform boundedness does not ensure a pointwise convergent subsequence}
    On \(S=[0,2\pi]\), consider again \(f_n(x)=\sin(nx)\). Suppose that a
    subsequence \(\{f_{n_k}\}_{k\in\naturalnumbers}\), with \(n_k\to\infty\),
    converges pointwise. Set
    \[
        g_k(x)=\bigl(f_{n_k}(x)-f_{n_{k+1}}(x)\bigr)^2.
    \]
    Then \(g_k\to0\) pointwise and \(0\leq g_k\leq4\). The dominated
    convergence theorem would therefore give
    \[
        \lim_{k\to\infty}\integral[0][2\pi][g_k(x)][x]=0.
    \]
    On the other hand, orthogonality of the sine functions gives, for every
    \(k\in\naturalnumbers\),
    \[
        \integral[0][2\pi][g_k(x)][x]
        =\integral[0][2\pi][\bigl(\sin(n_kx)-\sin(n_{k+1}x)\bigr)^2][x]
        =2\pi,
    \]
    a contradiction. Thus no such pointwise convergent subsequence exists.
\end{snippetexample}

\begin{snippetexercise}{sine-functions-orthogonality-exercise}{Orthogonality of sine functions}
    Prove that, for all \(i,j\in\naturalnumbers\),
    \[
        \integral[0][2\pi][\sin(ix)\sin(jx)][x]=
        \begin{cases}
            0 & i\neq j,\\
            \pi & i=j.
        \end{cases}
    \]
\end{snippetexercise}

\end{document}
